\section{Method}
\label{sec:method}

% SOURCE OF TRUTH: ../docs/PROTOCOL.md. That document is at Version 2 as this is
% written; Version 2 added two criteria on 2026-08-28 which are the subject of
% the ceiling post-mortem (Appendix~C). The seven-arm study ran under Version 2.

The evaluation reported here was conducted under Version 2 of the standing
protocol, adopted on 2026-08-28. Version 2 added two corpus-admission criteria
for reasons \Cref{sec:noreplication} summarises and \Cref{sec:ceiling}
reports in full. No rule under which this study ran was amended
retrospectively.

\subsection{Design}

Seven arms over the same items, paired by item. Only the system prompt differs.

\begin{center}
\begin{tabular}{ll}
\toprule
arm & system prompt \\
\midrule
\texttt{off} & task framing and response-format contract only \\
\texttt{on} & control plus the initial seed skill \\
\texttt{placebo} & control plus word- and structure-matched filler for \texttt{on} \\
\texttt{placebo-gepa} & control plus length- and structure-matched filler for GEPA's winner \\
\texttt{gepa} & control plus GEPA's frozen winner \\
\texttt{placebo-skillopt} & control plus length- and structure-matched filler for SkillOpt's winner \\
\texttt{skillopt} & control plus SkillOpt's frozen winner \\
\bottomrule
\end{tabular}
\end{center}

Three invariants hold across every arm.

\begin{enumerate}
  \item \textbf{The response-format contract is in every arm.} A skill that
    also teaches the model how to format an answer, tested against a control
    that was never told, measures instruction-following.
  \item \textbf{The placebos are matched on length and on structure.} To
    eliminate the length confound, each treatment arm is evaluated against its
    own matched placebo: \texttt{skillopt} against \texttt{placebo-skillopt},
    \texttt{gepa} against \texttt{placebo-gepa}, and \texttt{on} against
    \texttt{placebo}. \texttt{off} is reported alongside, outside the family,
    because it answers a different question: whether any document helps at all.
  \item \textbf{The option menu is held constant across arms.} Removing the
    explicit label menu moved trajectory accuracy by 14 to 40 points in
    AgentAtlas v1~\citep{agentatlas2026}. We cite v1 because v2 withdrew the
    number and replaced it with ``mapped label agreement can change
    substantially'', alongside a caution against reading its synthetic study
    as a model comparison. The figure therefore justifies holding the menu
    constant and nothing else; we do not compare it to our own effects.
\end{enumerate}

All seven arms are built by one function, \texttt{solvers.arms.build\_arm}, and
told apart in the records by the SHA-256 of the body each carries. An arm that
silently received the wrong body would be a different content hash in the
record.

\subsection{What the placebos control: eliminating the length confound}
\label{sec:length-confound}

In an earlier five-arm pilot, a single placebo matched only the seed skill,
leaving a length confound for the evolved winners. Here, each evolved winner is
paired with a length- and structure-matched placebo of its own, so that the
comparison isolates instructional content from prompt length.

Every record carries \texttt{input\_tokens}. The \texttt{off} arm sends the task
framing and the format contract and no document, so its prompt is the shared
prefix and each other arm's document is the difference. Per item that difference
has zero variance, so these are exact sizes.

\begin{center}
\begin{tabular}{lrr}
\toprule
arm & body tokens & $\times$ placebo \\
\midrule
\texttt{placebo} & \NUM{\bodyPlaceboTokens} & \NUM{\bodyPlaceboRatio} \\
\texttt{on} & \NUM{\bodyOnTokens} & \NUM{\bodyOnRatio} \\
\texttt{placebo-gepa} & \NUM{\bodyPlaceboGepaTokens} & \NUM{\bodyPlaceboGepaRatio} \\
\texttt{gepa} & \NUM{\bodyGepaTokens} & \NUM{\bodyGepaRatioVsPlaceboGepa} \\
\texttt{placebo-skillopt} & \NUM{\bodyPlaceboSkilloptTokens} & \NUM{\bodyPlaceboSkilloptRatio} \\
\texttt{skillopt} & \NUM{\bodySkilloptTokens} & \NUM{\bodySkilloptRatioVsPlaceboSkillopt} \\
\bottomrule
\end{tabular}
\end{center}

Each treatment arm is matched within 16\%--25\% by its dedicated placebo.
With matched placebos in place, any observed difference between an evolved
winner and its control cannot be attributed to prompt length or sectioning.

\subsection{Scope of the verdict}
\label{sec:arenas}

The project that ran this study registers every model against a tier that
decides whether a run on it may support a public claim about the shipped
skill, and the study's target, \NUM{\studyModel}, is registered at the
development tier. The result here is therefore published and controlled, and
scoped to one small model; it supports no claim about the skill at large.
\Cref{sec:noreplication} explains why a larger-model replication does not
exist.

\subsection{Pre-registration, verified rather than asserted}

The prediction for this study was committed to the repository's append-only
notebook before the first call. What makes that more than a claim is a gate: a
published run declares the notebook entry carrying its prediction, and the
provenance check refuses the run unless that entry's first commit is a git
ancestor of the commit the run was made at. Ancestry is not editable after the
fact without destroying the timestamps the method rests on.

The registered family was fixed before any data existed: three comparisons,
\texttt{on} vs \texttt{placebo}, \texttt{gepa} vs \texttt{placebo-gepa}, and
\texttt{skillopt} vs \texttt{placebo-skillopt}, one-sided in the direction of
the arm helping, applied to each item set separately. That is a family of three
per set and six registered comparisons in all, and the two sets are never
pooled. \texttt{off} sits outside the family by registration.

\subsection{Statistics: the pre-registered dual bar}

The pre-registration established a dual statistical bar:
\begin{enumerate}
  \item \textbf{Item level:} McNemar's exact test on discordant pairs,
    item-matched, one-sided, with Holm correction across the family of three
    within each set ($q < 0.05$).
  \item \textbf{Template level:} A permutation cluster sign-flip test across
    templates, requiring $p < 0.0167$. With \NUM{\templatesUnseen} unseen
    templates, the permutation support is $2^7 = 128$, providing a theoretical
    floor of $1/128 = 0.0078$ and sufficient power to detect authentic
    cross-template gains.
\end{enumerate}

Choosing the unit after seeing the data would be the defect this design exists
to avoid, so both are required for rejection. \Cref{sec:clustered} reports both
columns side by side. We report the discordant split beside every registered
$p$-value.

\subsection{Controls run before the results were read}

\paragraph{Two-pass replication and A/A.} The entire evaluation was executed
across two independent passes (\texttt{pass-1} and \texttt{pass-2}) to evaluate
reproducibility across all seven arms. In addition, the shared placebo was
scored an extra time over all \NUM{\totalItems} items into an A/A control pass
(14,700 calls in total). \Cref{sec:results} reports the outcome.

\paragraph{Falsifier.} A standing rule of the protocol, checked by hand
before any comparison was read: some possible response must score above zero
for every arm, and the scorer must read the same object in each. The rule
exists because an estimator that cannot return a non-zero value produces a
clean run and a plausible number, and nothing else about the run looks
wrong. \Cref{sec:results} reports the observation that satisfies it.

\paragraph{Residency.} Model residency is pinned at 60 minutes on every
request. A probe run hours before the study found two items answering
deterministically to whether the model had just been loaded, which without the
pin would have put a venue artifact inside the A/A control. The pin is a
property of the provider, and \Cref{sec:appendix-harness} states which other
settings share that status.
